\begin{longtable}{p{1.8cm} p{6.5cm} p{6.2cm}}
\caption{Cronograma por semanas, fases e marcos}
\label{tab:cronograma} \\
\toprule
\textbf{Semanas} & \textbf{Fase / Atividade} & \textbf{Critério de conclusão} \\
\midrule
\endfirsthead

\multicolumn{3}{c}{\tablename\ \thetable\ -- continuação} \\
\toprule
\textbf{Semanas} & \textbf{Fase / Atividade} & \textbf{Critério de conclusão} \\
\midrule
\endhead

\bottomrule
\multicolumn{3}{r}{\textit{continua}} \\
\endfoot

\bottomrule
\endlastfoot

1 & Análise exploratória da base; definição do protocolo de anotação & Estatísticas descritivas documentadas; protocolo escrito \\
2 & Geração de candidatos (blocagem produto--dicionário e produto--produto; supervisão distante) e configuração do treino & Pipeline de \textit{bootstrapping} funcional \\
3--5 & Anotação manual do conjunto de teste padrão-ouro (fatias produto--produto e produto--dicionário) e teste-reteste \newline \textit{Redação: Introdução} & Conjunto de teste anotado; Kappa intra-anotador calculado \\
6 & Partição anti-vazamento (nenhum produto em mais de uma partição) e aumentação sintética do treino \newline \textit{Redação: Revisão de Literatura} & Partições versionadas; conjunto de treino gerado \\
7--8 & Implementação e avaliação dos baselines clássicos \newline \textit{Redação: Metodologia} & Tabela de baselines fechada; harness reproduzível; checkpoint: $\geq$60\% do TCC redigido \\
9 & Configuração do SLM (dois adaptadores LoRA bf16), formatação dos pares, seeds e configs & Pipeline de treino executa de ponta a ponta \\
10--11 & Fine-tuning, ajuste de hiperparâmetros e mineração de negativos difíceis & Modelo treinado; curvas de treino registradas \\
12 & Avaliação do modelo vs. baselines, por fatia; teste de significância estatística & Tabela comparativa; hipótese central testada \\
13 & Ablação enxuta (aumentação; negativos difíceis); análise de pipeline progressivo & Dois estudos de ablação documentados; economia computacional do pipeline progressivo quantificada \\
14 & Análise de erro qualitativa por tipo de ruído & Taxonomia de erros e exemplos \\
15 & Redação de Resultados, Discussão e Conclusão; entrega de rascunho completo ao orientador no início da semana & Rascunho completo entregue; retorno do orientador obtido até o final da semana \\
16 & Incorporação do feedback do orientador, revisão final, formatação ABNT e preparação da defesa & TCC entregue; slides prontos \\

\end{longtable}
\vspace{-4mm}
\fonte{Elaborado pelo autor}
